\section{Discussion}
\label{sec:discussion}

The matrix results separate expressivity, approximation, and
preconditioning quality.  LMI feasibility determines whether a diagonal or
block family can reach a condition target.  Affine-invariant projection
determines the nearest structured geometry.  Condition-number minimization
selects the best preconditioner.  These operations agree in special regimes
and separate in general.

For Kronecker matrices, the factor scale gauge is part of the geometry rather
than a numerical nuisance.  Quotienting the gauge gives a complete totally
geodesic manifold, a closed-form intrinsic line element, and a unique AIRM
projection.  The partial-trace normal equations expose the projection
residual, while the extreme-state marginal condition determines whether that
projection is also condition-optimal.  The marginal-mismatch residual gives a
computable first-order measure of failure.  Every two-level relative spectrum
forces projection optimality.  Every \(2\times2\) Kronecker projection is also
condition-optimal, making the explicit \(2\times3\) separation
dimension-minimal.  Thus the universal guarantee does not extend to
unrestricted richer spectra in the next possible factor dimension.

At the tangent-space level, the same orthogonality calculation extends to a
product of \(k\) SPD factors.  In whitened coordinates its tangent log space
is the sum of the \(k\)
single-mode terms
\[
  I\otimes\cdots\otimes I\otimes X_j\otimes I\otimes\cdots\otimes I.
\]
Repeating the two-factor normal-space proof shows that the projection equations
require every one-mode partial trace of the log residual to vanish, and that
projection optimality is equivalent to choosing admissible density matrices
in the two extreme eigenspaces with matching marginals on every mode.
Fixed-basis dual formulas
and dimension-minimal separation for three or more factors require separate
analysis; the two-factor classification proved here does not presume either.

The exact penalty and proximal contraction provide a global variational
baseline; their rate is conditional on solving each proximal subproblem.
Likewise, the Armijo theorem gives a convergent projection solver but does not
claim that a dense Kronecker projection is inexpensive at neural-network
scale.  Practical implementations require matrix-free partial traces,
structured eigensolvers, and estimator-specific concentration bounds.

The strongest statements are deterministic finite-dimensional matrix
theorems.  Time-dependent curvature, damping, and optimizer states can be
inserted only after a concrete algorithm-to-geometry realization is specified.
The Loewner proxy results state precisely what must be supplied by such a
statistical layer.

The analytic certificates are exact-arithmetic statements.  The interval-safe
form remains rigorous when distance, residual, and condition number are
enclosed by validated numerical bounds and the derived comparisons are
evaluated with outward rounding.  Ordinary floating-point residuals
without such enclosures remain diagnostics; near multiple extreme eigenvalues,
the marginal condition should be solved over the full extreme eigenspaces
rather than from a single numerically selected eigenvector.

In summary, structured preconditioning has three distinct certificates:
reachability of a target, proximity to a full geometry, and
condition-optimality inside the family.  Affine-invariant matrix geometry
places all three in one framework and makes their coincidence and separation
mathematically explicit.
